\documentclass{beamer}
\usepackage[utf8]{inputenc}

\usetheme{Madrid}
\usecolortheme{default}
\usepackage{amsmath,amssymb,amsfonts,amsthm}
\usepackage{txfonts}
\usepackage{tkz-euclide}
\usepackage{listings}
\usepackage{adjustbox}
\usepackage{array}
\usepackage{tabularx}
\usepackage{gvv}
\usepackage{lmodern}
\usepackage{circuitikz}
\usepackage{tikz}
\usepackage{graphicx}

\setbeamertemplate{page number in head/foot}[totalframumber]

\usepackage{tcolorbox}
\tcbuselibrary{minted,breakable,xparse,skins}

\definecolor{bg}{gray}{0.95}
\lstset{
    language=C,
    basicstyle=\ttfamily\small,
    keywordstyle=\color{blue},
    stringstyle=\color{orange},
    commentstyle=\color{green!60!black},
    numbers=left,
    numberstyle=\tiny\color{gray},
    breaklines=true,
    showstringspaces=false,
}

%------------------------------------------------------------
\title{1.4.27}
\date{September 13,2025}
\author{EE25BTECH11011 - Navya Banavathu}

\begin{document}

\frame{\titlepage}

\begin{frame}{Question}
If $\vec{a}$ and $\vec{b}$ are the position vectors of $A$ and $B$ respectively, find the position vector of a point $C$ in $BA$ produced such that $BC$ = $1.5$ $BA$
\end{frame}

\begin{frame}{Solution}
Given the position vectors of $\vec{A}$ and  $\vec{B}$  are  $\vec{a}$ and  $\vec{b}$ \\
Let the position vector of point $\vec{C}$ be  $\vec{c}$ 

The vector {BC} is given by:
\begin{align}
\vec{BC}=   1.5{BA}\\
(\vec{c}-\vec{b}) = 1.5\, (\vec{b}-\vec{a})
\end{align}

we have,
\begin{align}
  \vec{AC}=\vec{BC}-\vec{BA}\\
  (\vec{c}-\vec{a}) = (\vec{c}-\vec{b}) - (\vec{a}-\vec{b}) \\
  (\vec{c}-\vec{a}) = (1.5 - 1)(\vec{a}-\vec{b}) \\
  (\vec{c}-\vec{a}) = 0.5\,(\vec{a}-\vec{b})
  \end{align}

  Thus,
\begin{align}
 (\vec{c}-\vec{a}) : (\vec{b}-\vec{c}) = 0.5 : 1.5 = 1 : 3 
\end{align}
Therefore, $ C $ divides $ AB $ externally in the ratio $ 1:3 $
\end{frame}

\begin{frame}{Section Formula(external division)}

Writing the ratio as $ k:1 $ 
we have 
\begin{align}
k=3
\end{align}
By the section formula (external division), the position vector of $C$  is
\begin{center}
\begin{align}
\vec{c}= \frac{k\vec{b} - \vec{a}}{k-1}, \quad (k=3).
\end{align}

\begin{align}
\vec{c} = \frac{3\vec{b} - \vec{a}}{2}.
\end{align}

\begin{align}
\therefore \quad \vec{C} =  \tfrac{3}{2}\vec{b} - \tfrac{1}{2}\vec{a} .
\end{align}

The above equation can be written as:
\begin{align}
\mathbf{c} = (\mathbf{a} \quad \mathbf{b})
\begin{pmatrix}
\frac{3}{2} \\
-\frac{1}{2}
\end{pmatrix}
\end{align}
\end{center}
\end{frame}


\begin{frame}[fragile]{ Code (C)}
\begin{block}{C Program}
\begin{lstlisting}[language=C, basicstyle=\ttfamily\small, keywordstyle=\color{blue}, frame=single]
#include <stdio.h>

int main() {
    float ax, ay, bx, by;
    printf("Enter coordinates of A (ax ay): ");
    scanf("%f %f", &ax, &ay);
    printf("Enter coordinates of B (bx by): ");
    scanf("%f %f", &bx, &by);
    int k = 3;
    float cx = (k*bx - ax) / (k - 1);
    float cy = (k*by - ay) / (k - 1);
    printf("Position vector of C is: (%.2f, %.2f)\n", cx, cy);
    return 0;
}
\end{lstlisting}
\end{block}
\end{frame}


\begin{frame}[fragile]{Python Code:Import and Define points}
\begin{block}{Part 1}
\begin{lstlisting}[language=Python, basicstyle=\ttfamily\small, keywordstyle=\color{blue}, frame=single]
import matplotlib.pyplot as plt

# --- Data ---
# Position vectors (points A and B)
A = (2, 3)   # Example coordinates of A
B = (5, -1)  # Example coordinates of B

# Ratio for external division (C divides AB externally in 3:1)
k = 3   
# Section formula for external division:
# C = (k*b - a) / (k - 1)
Cx = (k*B[0] - A[0]) / (k - 1)
Cy = (k*B[1] - A[1]) / (k - 1)
C = (Cx, Cy)
print(f"Position vector of C is: {C}")

\end{lstlisting}
\end{block}
\end{frame}

\begin{frame}[fragile]{Python Code: Plotting}
\begin{block}{Part 2}
\begin{lstlisting}[language=Python, basicstyle=\ttfamily\small, keywordstyle=\color{blue}, frame=single]
# --- Plotting ---
points = [A, B, C]
x_coords, y_coords = zip(*points)

plt.figure(figsize=(7, 6))

# Plot points
plt.scatter(x_coords, y_coords, color='red', s=80, zorder=5)

# Connect A-B (main line)
plt.plot([A[0], B[0]], [A[1], B[1]], color='blue', label='Line AB')

# Show extended line through A-B-C
plt.plot(x_coords, y_coords, '--g', label='Line with C')
\end{lstlisting}
\end{block}
\end{frame}


\begin{frame}[fragile]{Python Code: Finishing Plot}
\begin{block}{Part 3}
\begin{lstlisting}[language=Python, basicstyle=\ttfamily\small, keywordstyle=\color{blue}, frame=single]
# Labels
labels = ['A'+str(A), 'B'+str(B), 'C'+str(C)]
for (x, y), label in zip(points, labels):
    plt.text(x + 0.2, y + 0.2, label, fontsize=10)
plt.axhline(0, color='black', linewidth=0.5)
plt.axvline(0, color='black', linewidth=0.5)
plt.gca().set_aspect('equal', adjustable='box')
plt.xlabel('X-axis')
plt.ylabel('Y-axis')
plt.title('External Division (C divides AB in ratio 3:1)')
plt.legend()
plt.grid(True)
import os
# Your plotting code above...
plt.savefig("figs/fig1.png", dpi=300, bbox_inches="tight")
plt.show()
\end{lstlisting}
\end{block}
\end{frame}

\begin{frame}{Graphical Representation}
    \begin{figure}
        \centering
        \includegraphics[width=0.55\columnwidth]{figs/fig1.png}
        \caption{Point $C$ dividing $AB$ externally in ratio $3:1$}
    \end{figure}
\end{frame}


\end{document}